\section*{الفصل \ref{ch:g5:first-look-at-cells} --- نظرة أولى إلى الخلايا}

\begin{solution}{exo:g5:first-look-at-cells:1}
أصغر لَبِنة حية في \omterm{def:g1:living-or-not:living}{الكائن الحي}. وجسمك مصنوع من عشرات آلاف
المليارات منها.
\end{solution}

\begin{solution}{exo:g5:first-look-at-cells:2}
لأن \omterm{def:g5:first-look-at-cells:cell}{الخلايا} أصغر من أصغر ذرّة تستطيع \omterm{def:g1:the-five-senses:sense}{العين} رؤيتها. \omterm{def:g5:first-look-at-cells:microscope}{والمجهر} ---
وهو يكبّر الشرائح الرقيقة مئات المرات --- هو الذي كشفها.
\end{solution}

\begin{solution}{exo:g5:first-look-at-cells:3}
قشرة \omterm{ex:g4:plants-without-seeds:bulbtuber}{البصلة}: \omterm{def:g5:first-look-at-cells:cell}{خلايا} كالطوب في صفوف أنيقة، جدارًا إلى جدار. وكشط
الخدّ: \omterm{def:g5:first-look-at-cells:cell}{خلايا} أكثر استدارة متفرقة. والمشترك: أن المادتين مصنوعتان
كلتاهما من \omterm{def:g5:first-look-at-cells:cell}{خلايا}، تُظهر أكثرها بقعة داكنة صغيرة في داخلها.
\end{solution}

\begin{solution}{exo:g5:first-look-at-cells:4}
النواة --- وهي مركز قيادة \omterm{def:g5:first-look-at-cells:cell}{الخلية}.
\end{solution}

\begin{solution}{exo:g5:first-look-at-cells:5}
نحو $50$ \omterm{def:g5:first-look-at-cells:cell}{خلية} تمتد مليمترًا واحدًا؛ \omterm{def:g5:first-look-at-cells:cell}{والخلية} الواحدة نحو جزء من
خمسين من المليمتر --- أي $\qty{0.02}{mm}$.
\end{solution}

\begin{solution}{exo:g5:first-look-at-cells:6}
الازدياد طولًا: \omterm{def:g5:first-look-at-cells:cell}{خلايا} الجسم تنقسم وتتكاثر. والالتئام: \omterm{def:g5:first-look-at-cells:cell}{خلايا} تنقسم
لتعيد بناء الرقعة المصابة.
\end{solution}

\begin{solution}{exo:g5:first-look-at-cells:7}
واحدة. وهي \omterm{def:g4:animal-reproduction:sexual}{البويضة} المخصَّبة --- التحام \omterm{def:g4:animal-reproduction:sexual}{نطفة} واحدة \omterm{def:g4:animal-reproduction:sexual}{ببويضة} واحدة.
\end{solution}

\begin{solution}{exo:g5:first-look-at-cells:8}
المضاعفة ابتداءً من $1$ تعطي $2$ ثم $4$ ثم $8$ ثم $16$ ثم $32$:
فهناك $32$ \omterm{def:g5:first-look-at-cells:cell}{خلية} بعد $5$ انقسامات. وخمس مضاعفات أخرى تعطي $64$ ثم
$128$ ثم $256$ ثم $512$ ثم $1024$: أي $1024$ \omterm{def:g5:first-look-at-cells:cell}{خلية} بعد
$10$ انقسامات.
\end{solution}

\begin{solution}{exo:g5:first-look-at-cells:9}
مهما يختلفا حجمًا وشكلًا وحياةً، فكلاهما مبني من الصنف نفسه من
الوحدات الحية الصغيرة --- \omterm{def:g5:first-look-at-cells:cell}{الخلايا}. وخطة البناء المشتركة، وهي
مشتركة بين الكائنات الحية كلها، هي معنى ``\omterm{prop:g5:first-look-at-cells:principle}{وحدة الحياة}''.
\end{solution}

\begin{solution}{exo:g5:first-look-at-cells:10}
كلها: \omterm{def:g5:first-look-at-cells:cell}{فالخلية} المفردة تولد (من \omterm{def:g5:first-look-at-cells:cell}{خلية} أم تنقسم)، وتتغذى، وتنمو،
وتنقسم --- فتصنع كائنات حية جديدة مثلها --- وتموت. \omterm{def:g5:first-look-at-cells:cell}{فخلية} واحدة
تكفي لحياة كاملة.
\end{solution}

\begin{solution}{exo:g5:first-look-at-cells:11}
\omterm{def:g5:first-look-at-cells:cell}{خلايا} أكثر --- فحجم \omterm{def:g5:first-look-at-cells:cell}{الخلية} لا يكاد يتغير بين الفأر والفيل. والنمو
إذن ليس نفخًا \omterm{def:g5:first-look-at-cells:cell}{للخلايا} بل تكثيرًا لها: فالأجسام تكبر بانقسام بعد
انقسام، تمامًا كما تقول القضية.
\end{solution}
